% Delta-method error propagation: a non-linear g maps the input
% spread around mu_X to an output spread around g(mu_X), scaled by
% the local slope g'(mu_X).
\begin{figure}[htbp]
\centering
\begin{tikzpicture}[x=1cm, y=1cm, font=\small]

% Axes
\draw[->, thick] (-0.2, 0) -- (8.2, 0) node[right] {$x$};
\draw[->, thick] (0, -0.2) -- (0, 6.2) node[above] {$y = g(x)$};

% Nonlinear curve g(x) = 0.6*x^1.6 scaled
\draw[thick, engblue, smooth]
  plot[domain=0.3:7.6, samples=80] (\x, {0.30*pow(\x,1.55)});

% mu_X marker on x-axis
\def\muX{4.5}
\pgfmathsetmacro{\gmu}{0.30*pow(\muX,1.55)}
\draw[dashed, gray] (\muX, 0) -- (\muX, \gmu);
\draw[dashed, gray] (0, \gmu) -- (\muX, \gmu);
\node[anchor=north] at (\muX, -0.1) {$\mu_X$};
\node[anchor=east] at (-0.05, \gmu) {$g(\mu_X)$};

% Input spread band on x-axis
\draw[engred, very thick] (\muX-0.9, 0) -- (\muX+0.9, 0);
\node[anchor=north, engred, font=\scriptsize] at (\muX, -0.55)
  {input $\pm\sigma_X$};

% Tangent line at mu_X: slope g'(mu) = 0.30*1.55*mu^0.55
\pgfmathsetmacro{\slope}{0.30*1.55*pow(\muX,0.55)}
\draw[thick, engorange]
  (\muX-1.4, {\gmu - \slope*1.4}) -- (\muX+1.4, {\gmu + \slope*1.4});
\node[anchor=west, engorange, font=\scriptsize] at (\muX+1.0, {\gmu + \slope*1.2})
  {slope $g'(\mu_X)$};

% Output spread band on y-axis
\draw[engred, very thick]
  (0, {\gmu - \slope*0.9}) -- (0, {\gmu + \slope*0.9});
\node[anchor=east, engred, font=\scriptsize, rotate=90] at (-0.5, \gmu)
  {output $\pm g'(\mu_X)\sigma_X$};

\end{tikzpicture}
\caption[Delta-method error propagation]{The first-order delta
method. A non-linear function $g$ maps the input mean $\mu_X$ to the
output mean $g(\mu_X)$ and maps a small input spread $\sigma_X$ to
an output spread $|g'(\mu_X)|\,\sigma_X$, set by the local slope of
the tangent line. Where $g$ is steep the output uncertainty is
amplified; where $g$ is flat it is suppressed. The approximation
holds while $g$ is nearly linear across the input band; for a wide
band or a sharply curving $g$ the linearisation understates the true
spread and Monte Carlo propagation is the honest fallback. Source:
schematic by the authors.}
\label{fig:vol02:ch13:error-propagation}
\end{figure}
